% Chapter 6: RoPE Extrapolation Study
\chapter{RoPE Extrapolation in Streaming Diffusion Video Super-Resolution}
\label{ch:rope}

% Canonical sources: docs/notes/2026-07-16-rope-chapter-notes.md (numbers),
% docs/notes/2026-07-11-rope-extrapolation-findings.md (verdicts),
% reports/rope_timespace_sensitivity_matrix.md (standalone study incl. PI theory),
% reports/figures/{rope_probe_udm10_gt,dpp_causal_verdict}.md, results/rope_probe/*.

\section{Motivation and Research Questions}
\label{sec:rope:motivation}

The real-model evaluation of Chapter~\ref{ch:experiments} left a pointed question open:
FlashVSR --- within $0.012$ of the top composite score and the outright winner on
identity preservation --- exhibits the worst long-range semantic drift (sub-metric
D$''$) of the three methods --- and FlashVSR is precisely the method whose streaming
design lets its temporal RoPE index grow without bound as the video lengthens. Its
stock frequency table caps a single pass at 4{,}089 frames ($\approx 2.3$ minutes),
and its trained temporal window covers only 81 frames (21 latent frames after
$4\times$ VAE compression); its trained spatial extent is 48 latent rows (from the
$768\times1408$ distillation resolution). Position-encoding extrapolation is
therefore the natural suspect for long-video degradation, and the same suspicion
extends to the spatial axes at high resolution.

The study decomposes positional extrapolation into two falsifiable hypotheses:

\begin{description}
    \item[H0 (Magnitude).] Absolute position offsets beyond the trained window
        degrade quality.
    \item[H1 (Geometry).] Distortion of relative-distance geometry (dilation or
        compression of position spacing) degrades quality.
\end{description}

The decomposition matters because streaming inference stresses exactly the
magnitude axis (indices grow forever) while window extension stresses the geometry
axis (more positions inside one attention span). Which one costs quality determines
where engineering effort should go.

\section{The Position-Injection Instrument}
\label{sec:rope:instrument}

A runtime position-injection hook swaps any one of the three RoPE frequency tables
(t/h/w) --- a duck-typed replacement of the DiT's per-axis frequency-table
attribute --- for a wrapper that transforms baseline positions: shift
($p \mapsto p + k$), stretch ($p \mapsto s \cdot p$), continuous-position
interpolation, and modulo bounding, with on-demand table extension past the stock
1{,}024 rows. Continuous interpolation evaluates the RoPE formula directly at
fractional positions rather than interpolating table rows (averaging complex
phases would be incorrect); evaluated at integer positions it reproduces the stock
table's effect bit-for-bit, providing an identity condition at the numerics floor.

The instrument satisfies three faithfulness gates: (i) \emph{bit-exact when idle}
--- with the hook installed but no transform active, output drift against stock is
$0.0$ on a measured $0.0$ nondeterminism floor; (ii) \emph{engagement-checked} ---
every run asserts the transform actually fired; (iii) \emph{repository-pristine}
--- zero modification of the tagged FlashVSR codebase, with 48 unit tests covering
the hook. Content and compute are identical across conditions; only the position
labels change, which is what makes the comparisons causal.

\section{Experimental Design}
\label{sec:rope:design}

Seven experimental arms cover the two hypotheses per axis:
(1)~temporal shift sweeps up to position 996; (2)~temporal stretch sweeps,
integer and continuous, $s \in [0.5, 250]$; (3)~ground-truth scoring on
DOVE-UDM10 (10 clips $\times$ 29 frames, realistic degradation; PSNR/SSIM/LPIPS
via pyiqa in the DOVE RGB convention, on the protocol whose stock score ---
24.02~dB beside MGLD's verified 24.23 --- anchors every number); (4)~the D$''$
causal check on the 2{,}412--5{,}000-frame real-SR videos: stock segmented vs
true single-pass vs magnitude-bounded (positions cycled mod 336) arms;
(5)~spatial (h/w) label sweeps on UDM10; (6)~a real resolution-extrapolation
ladder on YouHQ40 (8 clips, $1080^2$ GT): latent grids
$48^2 \to 72^2 \to 80^2 \to 88^2 \to 96^2$, i.e.\ $1\times$ to $2\times$ the
trained spatial extent (input padding quantises grids to multiples of 8), stock
vs spatial-PI positions; and (7)~collapse-decomposition arms for the ladder's top
rung: content-correct scoring, sparsity pinned to the healthy $72^2$ value,
blur-isolation at fixed grid, and the intermediate grids.

A design principle throughout: \emph{score against real ground truth}. As
Section~\ref{sec:rope:selfconsistency} shows, self-consistency (similarity of
perturbed output to stock output) actively misleads --- a finding that
retroactively justifies the expense of GT-referenced scoring for every headline
claim.

\section{Results}
\label{sec:rope:results}

\subsection{Absolute-Position Offsets Are Free}
\label{sec:rope:offsets}

H0 is rejected on every axis. A temporal shift of $+996$ latent positions ---
$\approx 50\times$ beyond the trained window --- costs $+0.001$~dB against ground
truth; spatial offsets of 8 and 24 latent rows cost $\pm 0.003$~dB; and offsets
remain free at every stretch level tested (no interaction). Translation
invariance holds empirically on all three RoPE axes, as the relative-position
algebra of RoPE predicts but distilled, finite-precision practice did not
guarantee. Two engineering consequences follow: the streaming design's unbounded
index growth is harmless \emph{per se}, and the stock 4{,}089-frame single-pass
ceiling is an implementation artefact of the fixed table --- with the extended
table, a 5{,}009-frame single pass ran to completion at 11.6~GiB VRAM.

\subsection{Relative-Distance Geometry Is What Bites}
\label{sec:rope:geometry}

Table~\ref{tab:rope:stretch} sweeps position-label stretch per axis with content
and compute fixed. Dilation costs monotonically on the temporal axis
($-0.12$~dB at $s{=}1.25$ through $-0.95$ at $s{=}3$); strong compression costs
on every axis ($s{=}0.5$: $-0.66$ to $-1.00$~dB); and mild compression ---
$s{=}0.75$, the position-interpolation zone --- is \emph{free on every axis}
($+0.01$/$+0.02$/$-0.03$~dB). A temporal caveat calibrates the small $s \leq 2$
costs: under streaming, FlashVSR's attention span is $\sim$8 latents, so relative
distances remain inside the trained 21-latent range until $s \approx 2.6$ ---
those costs are within-window geometry effects, and $s{=}3$ is the first
genuinely out-of-window point.

\begin{table}[h]
\centering
\caption{$\Delta$PSNR (dB) vs stock under position-label stretch $s$, per RoPE
axis (DOVE-UDM10 vs ground truth; content and compute fixed).}
\label{tab:rope:stretch}
\begin{tabular}{lrrr}
\hline
\textbf{stretch $s$} & \textbf{time (t)} & \textbf{height (h)} & \textbf{width (w)} \\
\hline
0.5 (compression) & $-0.66$ & $-1.00$ & $-0.83$ \\
\textbf{0.75 (PI zone)} & $\mathbf{+0.01}$ & $\mathbf{+0.02}$ & $\mathbf{-0.03}$ \\
1.25 & $-0.12$ & $-0.05$ & $-0.02$ \\
1.5 & $-0.22$ & $-0.07$ & $-0.03$ \\
2.0 & $-0.25$ & $-0.15$ & $-0.07$ \\
3.0 (extreme) & $-0.95$ & $\mathbf{-2.55}$ & $\mathbf{-2.30}$ \\
\hline
\end{tabular}
\end{table}

\begin{figure}[h]
\centering
\includegraphics[width=0.9\textwidth]{time_vs_space_stretch.png}
\caption{Per-axis sensitivity to position-label stretch. The PI zone
($s{=}0.75$) is free on every axis; dilation costs grow monotonically, with the
spatial axes $\sim$2.5$\times$ more sensitive than time at extreme distortion.}
\label{fig:rope:stretch}
\end{figure}

At extreme geometry distortion the axes separate sharply: space is
$\sim$2.5$\times$ more sensitive than time ($-2.3$/$-2.6$~dB vs $-0.95$ at
$s{=}3$). The continuous-position builder that underlies these numbers was
validated by its identity condition (53.4~dB self-PSNR $=$ the numerics floor)
and revealed that genuine fractional-position interpolation is cheaper than
integer-rounded approximations by $+0.5$--$1.5$~dB --- integer rounding had been
silently inflating compression costs.

\subsection{Self-Consistency Does Not Predict Quality}
\label{sec:rope:selfconsistency}

Conditions $s{=}0.75$ and $s{=}1.5$ produce \emph{identical} output change
relative to stock (self-PSNR 31.8 for both) yet have opposite ground-truth
verdicts ($+0.01$ vs $-0.22$~dB). How much a perturbation changes the output
says nothing about whether the change damages quality. This is a methodological
finding for the field: position-perturbation studies that score only against the
unperturbed model's output --- the cheaper and common protocol --- cannot
distinguish benign from harmful interventions.

\subsection{Long-Video Drift Is Not Positional}
\label{sec:rope:drift}

The D$''$ causal check runs three arms on the full-length real-SR videos: stock
segmented inference, a true single pass (positions reaching latent 1{,}252 on the
extended table), and a magnitude-bounded arm (positions cycled mod 336). If
positional magnitude drove the long-video semantic drift measured in
Chapter~\ref{ch:experiments}, these arms would separate. They are
\textbf{indistinguishable --- under 1\% relative spread on every video}.
FlashVSR's weakest benchmark cell belongs to its streaming KV-cache/generation
mechanism, not to position encoding. The check also certifies retroactively that
the Chapter~\ref{ch:experiments} benchmark row was fair: segmented processing did
not disadvantage FlashVSR.

\subsection{Extreme Retrieval Distances}
\label{sec:rope:extreme}

Two questions posed by the research group --- ``test extreme stretches'' and
``what does it cost when a late frame retrieves an early one at its \emph{true}
large relative distance, rather than a re-based small one?'' --- reduce to the
same experiment. FlashVSR already uses true absolute indices (it never re-bases);
what keeps realised relative distances small is the $\sim$8-latent KV cache,
which evicts old keys long before a distance-1{,}000 attention edge could form.
(The re-based, window-local alternative is exactly SeedVR2's design --- the
planned round-2 contrast model.) Holding content fixed and inflating only the
position label isolates what that latent danger would cost
(Table~\ref{tab:rope:extreme}).

\begin{table}[h]
\centering
\caption{Temporal extrapolation at extreme effective distances (DOVE-UDM10 vs
ground truth, continuous positions; the stock table is bypassed by the
continuous builder).}
\label{tab:rope:extreme}
\begin{tabular}{lllr}
\hline
\textbf{stretch $s$} & \textbf{eff.\ distance} & \textbf{vs window ($\sim$20)} & \textbf{$\Delta$PSNR} \\
\hline
2 & $\sim$16 & window edge & $-0.25$ \\
5 & $\sim$40 & $2\times$ & $-1.34$ \\
10 & $\sim$80 & $4\times$ & $-1.03$ \\
20 & $\sim$160 & $8\times$ & $-0.31$ \\
50 & $\sim$400 & $20\times$ & $-1.57$ \\
125 & $\sim$1000 & $50\times$ (``frame 0 from frame 1000'') & $-1.33$ \\
250 & $\sim$2000 & $100\times$ & $-1.56$ \\
\hline
\end{tabular}
\end{table}

\begin{figure}[h]
\centering
\includegraphics[width=0.85\textwidth]{extreme_distance_temporal.png}
\caption{Temporal retrieval-distance dose-response: bounded, graceful, and
non-monotonic. Even 100$\times$ the trained window plateaus at $\sim$$-1.5$~dB.}
\label{fig:rope:extreme}
\end{figure}

Three readings. First, \textbf{the loss is bounded and saturates --- no
collapse}: even at 50--100$\times$ the trained window, quality sits at $-1.3$ to
$-1.6$~dB, barely worse than $s{=}3$'s $-0.95$; temporal extrapolation degrades
gracefully, the opposite of the spatial axes' steep $-2.3$/$-2.6$~dB at $s{=}3$
alone. Second, \textbf{the curve is non-monotonic}: $s{=}20$ recovers to
$-0.31$~dB between $\sim$$-1$~dB neighbours, because rotary phases are periodic
and certain large stretches accidentally re-align fast frequencies near
trained-like configurations --- \emph{where the phases land}, not how far they
moved, sets the damage. Third, the group's answer: the true-index design costs a
bounded $\sim$1.3~dB at distance 1{,}000 --- and FlashVSR never pays even that,
because its short cache prevents the far edge from forming. Absolute index plus
short cache is safe by construction, which is also why long-video drift is not
positional.

\subsection{Real Resolution Extrapolation}
\label{sec:rope:resolution}

The label sweeps above manipulate positions at fixed content; the resolution
ladder extrapolates for real, growing the latent grid from the trained extent
($48^2$, output $720^2$) to double it ($96^2$, output $1536^2$) on YouHQ40
(Table~\ref{tab:rope:ladder}).

\begin{table}[h]
\centering
\caption{YouHQ40 resolution ladder: PSNR vs ground truth per latent grid, stock
positions vs spatial-PI.}
\label{tab:rope:ladder}
\begin{tabular}{lllrr}
\hline
\textbf{output} & \textbf{latent grid} & \textbf{vs trained} & \textbf{stock} & \textbf{spatial-PI} \\
\hline
$720^2$ & $48^2$ & $1\times$ (trained) & 24.70 & $\equiv$ stock \\
$1152^2$ & $72^2$ & $1.5\times$ & 24.55 & 24.00 \\
$1280^2$ & $80^2$ & $1.67\times$ & 24.66 & --- \\
$1408^2$ & $88^2$ & $1.83\times$ & 24.71 & --- \\
$1536^2$ & $96^2$ & $2\times$ & \textbf{24.78} & --- \\
\hline
\end{tabular}
\end{table}

The ladder is \textbf{flat}: real resolution extrapolation is quality-free to at
least $2\times$ the trained spatial extent with stock positions --- adaptive
attention sparsity included (pinning the sparsity ratio to the healthy $72^2$
value changes the top rung by only $-0.03$~dB). The only positional effect on
the ladder is that \textbf{spatial-PI at $1.5\times$ hurts} ($-0.56$~dB versus
simply extrapolating): compressing positions by factor 0.67 lands in the
compression-cost zone of Table~\ref{tab:rope:stretch}. Spatially, compression
--- not extrapolation --- is the risky direction.

\paragraph{A methods note, recorded honestly.} An early scoring pass reported an
$-11$~dB ``collapse'' at the top rung. Decomposition arms cleared every
model-side suspect (sparsity pinning $\Delta$$-0.03$; intermediate grids flat;
no blur penalty from the upsampled input), tracing the number to a
scoring-geometry artefact: the padded output frame (content occupying 93.75\%)
had been resized against ground truth at mismatched magnification. This is the
second instance in this thesis where a scoring-path subtlety flipped a
conclusion (cf.\ Section~\ref{sec:rope:selfconsistency}); both were caught
because raw frames and re-scorable outputs were retained --- a reproducibility
practice worth stating as such.

\subsection{Why Position Interpolation Works --- and Where It Stops}
\label{sec:rope:pitheory}

A fair objection to the PI results: fractional positions produce rotary phases
the model never saw --- are those not out-of-distribution too? The resolution
\cite{chen_pi_2023} is the distinction between \emph{interpolation within} the
trained range and \emph{extrapolation beyond} it. Attention consumes relative
phase $(p - q)\,\theta_k$; training taught the model this function at roughly 20
discrete distances per frequency. PI places new inputs \emph{between} trained
points, where a learned smooth function is well-behaved; dilation places them
\emph{beyond} the trained interval, where the same function is unconstrained.

\begin{figure}[h]
\centering
\includegraphics[width=0.9\textwidth]{fig2_kernel.png}
\caption{The learned relative-phase kernel: PI queries it between trained
points (interpolation); dilation queries it beyond them (extrapolation).}
\label{fig:rope:kernel}
\end{figure}

The measurements above are this theory's per-axis fingerprint --- apparently the
first in a video diffusion model --- sharpened by a symmetry argument: the
rotation-change magnitude of a stretch depends only on $|s - 1| \cdot d$, which
is \emph{symmetric} between compression and dilation, yet measured quality is
strongly asymmetric ($s{=}0.75$ free everywhere; extrapolation up to
$-2.6$~dB). \emph{Where} phases land decides damage, not how far they moved.
The theory also predicts its own boundary: at $s{=}0.5$, neighbouring tokens sit
at relative distance $0.5$ --- below the minimum nonzero trained distance
($1.0$), inside the qualitatively different self-vs-neighbour transition ---
which is why the zero-shot PI-free zone is narrow ($\sim$0.75) and why the LLM
literature pairs aggressive PI with fine-tuning.

\section{Implications}
\label{sec:rope:implications}

\textbf{Position encoding is exonerated.} Offsets are free on every axis at any
magnitude tested; realistic geometry costs are small; the two failure modes that
motivated the study --- long-video drift and high-resolution degradation --- are
both demonstrably non-positional. Long-video improvement effort should target
the streaming KV-cache/generation mechanism; the resolution path is already free
to $2\times$.

\textbf{Guidance for window extension.} Temporal extension to $\sim$$1.33\times$
(21 $\to$ 28 latents) with continuous-PI positions is predicted RoPE-loss-free;
beyond that, temporal geometry costs grow but stay bounded and graceful --- even
$100\times$ the window plateaus at $\sim$$-1.5$~dB, so RoPE will not be the wall
a window extension hits. The validated instrument transfers directly: comparing
stock vs PI positions in extended-window runs isolates RoPE's causal share.

\textbf{Do not transplant LLM context-extension recipes blindly.} PI helps only
in its narrow zero-shot zone and actively hurts at $\leq$$1.5\times$ spatial
extension, where stock extrapolation is cheaper. Per-axis measurement, not
one-dimensional intuition, is the transferable practice.

\section{Limitations}
\label{sec:rope:limitations}

All ground-truth evidence is temporally within-window (29--31-frame GT clips);
deep-length ($>$81-frame) quality-vs-position curves lack ground truth, as
re-acquisition of long HR sources is blocked platform-side --- the path forward
is lab-provided long HR footage, for which conversion tooling is ready. The
study covers a single model (FlashVSR 1.3B distilled); SeedVR2, with
window-local positions, is the designated contrast for a second round. The
spatial trained extent is inferred from the distillation resolution rather than
training logs.
